\section{Conclusion}\label{sec:conclusion}

We have established a theoretical framework connecting classical proof
complexity measures to the convergence behavior of chain-of-thought
verification in large language models.  Our main results show that: (1)~the
number of CoT steps to convergence scales as $O(\log(1/\varepsilon)/\lambda)$,
where the convergence rate~$\lambda$ depends inversely on a composite proof
complexity measure (Theorem~\ref{thm:convergence}); (2)~asymptotic verification
accuracy decays exponentially in proof length under step-independence
(Theorem~\ref{thm:error-accumulation}); and (3)~self-consistency amplification
requires $\Theta(1/(p - 1/2)^2)$ chains, where the base accuracy~$p$ itself
depends on proof complexity (Theorem~\ref{thm:self-consistency}).  Computational
experiments on~129 synthetic instances across four formula families strongly
support these predictions, and a phase transition in verification difficulty is
identified at a critical complexity threshold~$C^* \approx 12.95$
(Proposition~\ref{prop:phase-transition}).

The central takeaway is that proof complexity theory provides the right lens
for understanding the limits of neural proof verification.  The framework
offers a principled basis for allocating test-time compute in mathematical
reasoning systems, replacing model-dependent difficulty estimates with
formula-intrinsic complexity measures.

Two immediate directions for future work are: (i)~empirical validation of the
predicted complexity--convergence relationships using state-of-the-art LLMs,
which would test Conjectures~\ref{conj:empirical} and~\ref{conj:space}; and
(ii)~extension of the framework beyond propositional resolution to first-order
and natural-language proofs, leveraging the Curry--Howard bridge introduced by
Perrier~\cite{perrier2025typed}.
